% !TeX program = xelatex
\documentclass[11pt,oneside]{article}
\usepackage{ProblemsetStyle}

\hypersetup{
    pdftitle={20260819 Problem set - Solutions},
    pdfauthor={김정우},
    pdfsubject={Science Quiz mathematics practice problem set solutions},
    pdfcreator={XeLaTeX}
}

\begin{document}

\problemsettitle{20260819 Problem set}

% Source check: International Mathematical Union, Fields Medals 2026.
\begin{problemsolution}{1}
International Mathematical Union의 2026 Fields Medal 공식 발표에 따르면, Hong Wang은 조화해석과 기하측도론에서의 연구와 함께 Fourier restriction, Falconer distance sets, 평면의 Furstenberg sets, 그리고 3차원 Kakeya problem에 대한 주요 진전으로 필즈상을 받았다.

따라서 정답은
\[
    \boxed{\text{Hong Wang}}
\]
이다.
\end{problemsolution}

\begin{problemsolution}{2}
적분을
\[
    I\coloneqq \int_0^1
    \frac{x^{2026}}{x^{2026}+(1-x)^{2026}}
    \,\dd x
\]
라 하자. $x\mapsto1-x$를 대입하면
\[
    I=\int_0^1
    \frac{(1-x)^{2026}}{(1-x)^{2026}+x^{2026}}
    \,\dd x.
\]
두 식을 더하면
\[
    2I
    =\int_0^1 1\,\dd x
    =1.
\]
따라서
\[
    \boxed{I=\frac12}.
\]
\end{problemsolution}

\begin{problemsolution}{3}
$\mathbf{A}^2=\mathbf{A}$이므로 $\mathbf{A}$는 image에서는 identity로, kernel에서는 $0$으로 작용한다. Rank가 $3$이므로 image는 3차원이고, rank--nullity theorem에 의해 kernel은 1차원이다. 따라서 $\mathbf{I}_4+\mathbf{A}$는 image에서 $2$배, kernel에서 $1$배로 작용한다. 그러므로
\[
    \det\!\left(\mathbf{I}_4+\mathbf{A}\right)
    =2^3\cdot1
    =\boxed{8}.
\]
\end{problemsolution}

% Source check: The Abel Prize, 2026 laureate announcement and citation.
\begin{problemsolution}{4}
2026 Abel Prize의 공식 수상 사유는 산술기하학에 강력한 도구를 도입하고 Mordell과 Lang의 오래된 디오판토스 추측들을 해결한 공로이며, 해당 수상자는 1986년 필즈상 수상자이기도 하다.

따라서 정답은
\[
    \boxed{\text{Gerd Faltings}}
\]
이다.
\end{problemsolution}

\begin{problemsolution}{5}
정사영 공식에 의해
\[
    \proj_{\mathbf{v}}\mathbf{x}
    =
    \frac{\langle\mathbf{x},\mathbf{v}\rangle}
    {\langle\mathbf{v},\mathbf{v}\rangle}\mathbf{v}.
\]
여기서
\[
    \langle\mathbf{x},\mathbf{v}\rangle
    =3\cdot1+1\cdot2+(-1)\cdot2
    =3
\]
이고
\[
    \langle\mathbf{v},\mathbf{v}\rangle
    =1^2+2^2+2^2
    =9.
\]
따라서
\[
    \proj_{\mathbf{v}}\mathbf{x}
    =\frac13(1,2,2)
    =
    \boxed{\left(\frac13,\frac23,\frac23\right)}.
\]
\end{problemsolution}

\begin{problemsolution}{6}
$2027$은 소수이고 $3,5,2026$은 모두 $2027$의 배수가 아니므로 Fermat's little theorem에 의해
\[
    3^{2026}\equiv1\pmod{2027},
    \qquad
    5^{2026}\equiv1\pmod{2027}.
\]
또한
\[
    2026\equiv-1\pmod{2027}
\]
이므로
\[
    2026^{2026}\equiv(-1)^{2026}=1\pmod{2027}.
\]
따라서
\[
    3^{2026}+5^{2026}+2026^{2026}
    \equiv3\pmod{2027}.
\]
그러므로 나머지는
\[
    \boxed{3}
\]
이다.
\end{problemsolution}

% Source check: Shuhong Gao, Counterexamples to the Jacobian conjecture in dimensions greater than two, arXiv:2608.00222.
\begin{problemsolution}{7}
2026년 7월에 $\C^3$에서 Jacobian conjecture의 반례가 공개되었고, 이어 공개된 Shuhong Gao의 프리프린트는 이 구성 원리를 일반화하여 모든 차원 $n>2$에서 반례를 제시한다. 반면 복소수체 위의 고전적인 $n=2$ 경우는 2026년 8월 19일 현재 여전히 미해결이다.

따라서 정답은
\[
    \boxed{(3)}
\]
이다.
\end{problemsolution}

\begin{problemsolution}{8}
적어도 하나의 주사위가 $6$인 순서 있는 결과는 모두 11개이다. 구체적으로
\[
    \begin{aligned}
        &(6,1),(6,2),\ldots,(6,6),\\
        &(1,6),(2,6),\ldots,(5,6)
    \end{aligned}
\]
이다. 이 중 합이 $10$ 이상인 것은
\[
    (6,4),(6,5),(6,6),(4,6),(5,6)
\]
의 5개이다. 따라서 조건부확률은
\[
    \boxed{\frac5{11}}.
\]
\end{problemsolution}

\begin{problemsolution}{9}
$A$에서 $B$로 가는 전체 함수의 수는 $3^4$이다. 전사가 되려면 $B$의 세 원소가 모두 함수값으로 나타나야 하므로 포함배제 원리를 적용하면
\[
    3^4
    -\binom31 2^4
    +\binom32 1^4
    =81-48+3
    =36.
\]
따라서 전사함수의 개수는
\[
    \boxed{36}
\]
이다.
\end{problemsolution}

% Source check: official IMO 2026 individual results.
\begin{problemsolution}{10}
2026 IMO 공식 개인 결과에서 대한민국 대표 Hyeonjun Lee는 여섯 문제 모두 7점을 받아 총점 42점으로 개인 공동 1위를 기록했다.

따라서 정답은
\[
    \boxed{\text{Hyeonjun Lee}}
\]
이다.
\end{problemsolution}

\end{document}
